% FreeCAD Cloud Agent Environment - Scoping Document
% Minimalist, Apple-inspired styling. Compile with: pdflatex environment-setup-scope.tex
\documentclass[11pt]{article}

% --- Page geometry: generous whitespace -------------------------------------
\usepackage[a4paper,top=28mm,bottom=26mm,left=30mm,right=30mm]{geometry}

% --- Typography: clean sans-serif throughout --------------------------------
\usepackage[T1]{fontenc}
\usepackage[utf8]{inputenc}
\usepackage[scaled=1.02]{helvet}
\renewcommand{\familydefault}{\sfdefault}
\usepackage{microtype}
\usepackage[document]{ragged2e}
\usepackage{setspace}
\setstretch{1.18}

% --- Muted, restrained color palette ----------------------------------------
\usepackage{xcolor}
\definecolor{appleink}{HTML}{1D1D1F}   % near-black body text
\definecolor{applegray}{HTML}{6E6E73}  % secondary text
\definecolor{appleblue}{HTML}{0071E3}  % single accent
\definecolor{hairline}{HTML}{D2D2D7}   % thin rules
\definecolor{panel}{HTML}{F5F5F7}      % soft panel fill
\color{appleink}

% --- Section styling: light, spacious, unnumbered ---------------------------
\usepackage{titlesec}
\setcounter{secnumdepth}{0}
\titleformat{\section}
  {\LARGE\bfseries\color{appleink}}{}{0em}{}
\titlespacing*{\section}{0pt}{2.4em}{0.8em}
\titleformat{\subsection}
  {\large\bfseries\color{appleink}}{}{0em}{}
\titlespacing*{\subsection}{0pt}{1.4em}{0.4em}

% --- Lists: tight and tidy --------------------------------------------------
\usepackage{enumitem}
\setlist{leftmargin=1.3em,itemsep=3pt,topsep=4pt,parsep=0pt}
\setlist[itemize]{label=\textcolor{appleblue}{\small\raisebox{0.15ex}{\textbullet}}}

% --- Tables -----------------------------------------------------------------
\usepackage{booktabs}
\usepackage{tabularx}
\usepackage{array}
\renewcommand{\arraystretch}{1.35}
\newcolumntype{L}{>{\RaggedRight\arraybackslash}X}

% --- Boxes for the summary panel --------------------------------------------
\usepackage[most]{tcolorbox}

% --- Links ------------------------------------------------------------------
\usepackage[hidelinks]{hyperref}
\hypersetup{colorlinks=true,urlcolor=appleblue,linkcolor=appleblue}

% --- Headers / footers: minimal ---------------------------------------------
\usepackage{fancyhdr}
\pagestyle{fancy}
\fancyhf{}
\renewcommand{\headrulewidth}{0pt}
\renewcommand{\footrulewidth}{0pt}
\fancyfoot[C]{\footnotesize\color{applegray}FreeCAD\quad\textbullet\quad Cloud Agent Environment\quad\textbullet\quad\thepage}

% Thin hairline helper
\newcommand{\hairrule}{\vspace{2pt}{\color{hairline}\hrule height 0.6pt}\vspace{2pt}}

% Eyebrow label helper (small, letter-spaced, uppercase)
\usepackage{soul}
\newcommand{\eyebrow}[1]{{\footnotesize\color{appleblue}\bfseries\so{\MakeUppercase{#1}}}}

\begin{document}
\thispagestyle{fancy}

% ============================ TITLE BLOCK ===================================
\begingroup
\setstretch{1.0}
\noindent\eyebrow{Cloud Agent Environment}\\[1.1em]
{\fontsize{30}{34}\selectfont\bfseries Development Environment\\ Setup \textemdash\ Scoping Document\par}
\vspace{0.7em}
{\large\color{applegray} Reproducible, prebuilt FreeCAD toolchain for autonomous Cloud Agents.\par}
\vspace{1.1em}
\noindent{\footnotesize\color{applegray} Repository \; \textbf{brucecodes2023/FreeCAD\_BT} \hfill Package manager \; \textbf{pixi (conda-forge)} \hfill Target \; \textbf{linux-64}}
\endgroup

\vspace{0.9em}
\hairrule
\vspace{1.4em}

% ============================ SUMMARY PANEL =================================
\begin{tcolorbox}[
  enhanced, colback=panel, colframe=panel, boxrule=0pt,
  arc=6pt, left=16pt, right=16pt, top=12pt, bottom=12pt]
{\eyebrow{Summary}}\\[0.6em]
FreeCAD is a large C++ / Python parametric CAD application built on OpenCASCADE,
Coin3D, and Qt6. This document scopes the work required to make the repository
\emph{fully usable} inside Cursor Cloud Agents: a deterministic base image, an
idempotent install phase that compiles the application from source, and an
end-to-end validation flow that proves the environment works before it is saved
as a reusable prebuilt snapshot.
\end{tcolorbox}

% ============================ OBJECTIVES ====================================
\section{Objectives}
\begin{itemize}
  \item Provide a one-command, reproducible setup that any future Cloud Agent inherits automatically.
  \item Reuse the project's canonical, pinned toolchain (\texttt{pixi} + \texttt{pixi.lock}) rather than inventing a parallel dependency story.
  \item Compile FreeCAD once during environment \texttt{install} so prebuilt snapshots boot ready-to-run.
  \item Demonstrate the environment end to end: build, automated tests, and a real modeling action.
\end{itemize}

% ============================ CURRENT STATE =================================
\section{Current State}
The repository already defines a complete, cross-platform build description. Key findings from the codebase analysis:

\begin{itemize}
  \item \textbf{Dependency management:} \texttt{pixi.toml} + \texttt{pixi.lock} pin every dependency from \texttt{conda-forge} (OpenCASCADE~7.8, Qt6~6.8, PySide6, Coin/pivy, SMESH, VTK, compilers, \texttt{ccache}, \texttt{mold}).
  \item \textbf{Build system:} CMake with presets in \texttt{CMakePresets.json}; \texttt{pixi} tasks wrap configure / build / install / test for each platform.
  \item \textbf{Reference CI:} \texttt{.github/workflows/sub\_buildPixi.yml} encodes the canonical Linux flow \textemdash{} \texttt{configure-release} \textrightarrow{} \texttt{build-release} \textrightarrow{} \texttt{install-release} \textrightarrow{} tests.
  \item \textbf{Submodules:} five git submodules (OndselSolver, GSL, Coin, pivy, AddonManager) fetched via the \texttt{initialize} task.
  \item \textbf{No prior Cloud Agent config:} no committed \texttt{.cursor/environment.json} existed before this work.
\end{itemize}

% ============================ SCOPE =========================================
\section{Scope}

\subsection{In scope}
\begin{itemize}
  \item A committed \texttt{.cursor/environment.json} plus an idempotent \texttt{.cursor/install.sh}.
  \item Installing \texttt{pixi} onto \texttt{PATH}, fetching submodules, and materializing the locked environment.
  \item Configuring, building, and installing FreeCAD with the \texttt{release} preset (matching CI).
  \item Local validation on the VM: full build, C++ / Python self-tests, and a headless modeling action.
  \item Snapshot, prebuilt-build verification, and a proposed configuration for review.
\end{itemize}

\subsection{Out of scope}
\begin{itemize}
  \item Changes to FreeCAD application or build-system source (environment configuration only).
  \item Non-Linux targets (macOS / Windows) \textemdash{} already covered by the project's own CI matrix.
  \item Dependency upgrades or \texttt{pixi.lock} regeneration.
  \item Packaging, release artifacts, or deployment.
\end{itemize}

% ============================ APPROACH ======================================
\section{Technical Approach}

\subsection{Base image}
Use Cursor's default Ubuntu image. It already provides \texttt{git}, \texttt{curl}, and
passwordless \texttt{sudo}; the full compiler toolchain arrives through \texttt{pixi}, so a
custom Dockerfile is unnecessary. The \texttt{pixi} binary is installed once to
\texttt{/usr/local/bin} so every agent shell finds it without shell-profile edits.

\subsection{Lifecycle placement}
\begin{itemize}
  \item \textbf{\texttt{install}} \textemdash{} durable, idempotent, source-derived work: install \texttt{pixi}, fetch submodules, resolve the locked environment, then configure / build / install FreeCAD. \texttt{ccache} keeps re-runs incremental.
  \item \textbf{\texttt{start} / \texttt{terminals}} \textemdash{} none required. FreeCAD is a desktop application, not a long-running service; GUI runs happen on demand under a virtual display (\texttt{xvfb}).
\end{itemize}

\subsection{Prebuilt snapshots}
Because the compile happens in \texttt{install}, the resulting binaries are baked into the
environment snapshot. Fresh agents boot with FreeCAD already built and do not pay the
compilation cost again.

% ============================ DELIVERABLES ==================================
\section{Deliverables}
\begin{itemize}
  \item \texttt{.cursor/environment.json} \textemdash{} declares the environment and its install entry point.
  \item \texttt{.cursor/install.sh} \textemdash{} idempotent bootstrap + build script.
  \item This scoping document (source + compiled PDF).
  \item Validation evidence: build logs, passing tests, and a headless modeling artifact.
\end{itemize}

% ============================ VALIDATION ====================================
\section{Validation Plan}
\noindent\begin{tabularx}{\linewidth}{@{}p{0.32\linewidth}L@{}}
\toprule
\textbf{Check} & \textbf{Success criterion} \\
\midrule
Configure & \texttt{configure-release} resolves all dependencies without error. \\
Build & \texttt{build-release} compiles every target (8092) cleanly. \\
Install & \texttt{install-release} stages the runnable application. \\
CLI self-tests & \texttt{FreeCADCmd -t 0} passes the test suite. \\
C++ unit tests & \texttt{ctest} green across module test runners. \\
Modeling action & A scripted part is created, and geometric properties (volume / area) verified. \\
Idempotence & Re-running \texttt{install} converges with no rebuild churn. \\
Fresh agent & An agent booted from the prebuilt build repeats the checks. \\
\bottomrule
\end{tabularx}

% ============================ RISKS =========================================
\section{Risks \& Mitigations}
\noindent\begin{tabularx}{\linewidth}{@{}p{0.40\linewidth}L@{}}
\toprule
\textbf{Risk} & \textbf{Mitigation} \\
\midrule
Long initial compile (thousands of TUs) & Runs once in \texttt{install}; snapshot bakes in binaries; \texttt{ccache} accelerates re-runs. \\
Headless GUI (no physical display) & Exercise the GUI/tests under \texttt{xvfb}; core validation via the CLI path. \\
Network access to conda-forge / GitHub & Uses the pinned \texttt{pixi.lock}; egress is currently unrestricted for the required hosts. \\
\bottomrule
\end{tabularx}

% ============================ HANDOFF =======================================
\section{Handoff}
Once local validation and prebuilt-build verification pass, the configuration is
proposed for review. Saving it in the Environment panel is the single remaining
manual action and makes the tested setup available to all future Cloud Agents.

\end{document}
